\section{Codebooks}
\label{apx-sec-codebook}

\subsection{Design Concepts}
\label{apx-design-concepts}

\begin{table*}[!ht]
\caption{Codebook of design concepts (DC1--DC6) used to organise workshop outputs in Layer~1.}
\Description{A codebook table listing the six design concepts (DC1–DC6) used to organise workshop outputs in Layer 1, including concept labels and brief descriptions defining each concept’s focus.}
\label{tab:codebook-design-concepts}
\begin{tabular}{@{} p{0.08\linewidth} p{0.22\linewidth} p{0.64\linewidth} @{}}
\toprule
\textbf{ID} & \textbf{Concept label} & \textbf{Brief description} \\
\midrule
DC1 & Tactile Stroll / Atrium Path &
A break-route composed of varied underfoot terrain---uneven flooring, trunks, stepping elements, height differences, and playful obstacles---that induces attentional shift and mental reset through kinaesthetic engagement. \\
DC2 & Diversified Entrance &
A multi-path entrance zone using differentiated stones, tactility, and soft-edged ``desire path'' to support intuitive wayfinding and subtle, ground-level sensory variation. \\
DC3 & Textural Office &
An open office environment structured through moving semi-transparent curtains, operable façade elements, airflow modulation, and visible infrastructural textures that introduce gentle movement and light variation. \\
DC4 & In-Between-Headspace Space &
A transitional ``in-between'' space for mental reset, characterised by layered changes in temperature, sound, smell, and light, combined with a controlled sense of spatial randomness inspired by outdoor environmental variability. \\
DC5 & Biophilic Sound Study Nook &
A multi-purpose study nook where plants, soil, water, and  acoustic surfaces soften the soundscape to create an immersive space for relaxation, focused work, and socialisation. \\
DC6 & Temporal Staircase &
A staircase environment incorporating flowing water, stepping-stone logics, plants, and temporally modulated sound to reduce echoes, mark transitions, and introduce subtle acoustic and temporal variation. \\
\bottomrule
\end{tabular}
\end{table*}


\subsection{Layer 1 Design Space Concepts}
\begin{table*}[!ht]
\caption{Layer~1 codebook developed during content-based analysis and mapping the biophilic design space (Part 1 of 2. Table 2 follows.).}
\Description{Codebook for Layer 1 showing dimensions, codes, and definitions for Sensory Modality, Spatial Context, and Interaction Type.}
\label{tab:layer1-codebook}
\centering
\begin{tabular}{@{}
  p{0.18\textwidth}
  p{0.20\textwidth}
  p{0.58\textwidth}
@{}}
\toprule
\textbf{Dimension} & \textbf{Code} & \textbf{Definition} \\
\midrule
\multirow{6}{=}{\centering\textbf{Sensory Modality}} 
& Sight & How the eyes register light, movement, depth, colour, patterns. \\
& Sound & How the ears perceive acoustic layers, direction, texture, and shifts. \\
& Smell & How the nose detects atmospheric qualities, freshness, humidity, odours. \\
& Taste & How the mouth encounters flavour, edible elements, or atmospheric traces carried indoors. \\
& Touch & How the skin senses temperature, texture, and material contact. \\
& Kinesthetic & How the body perceives movement, balance, gait, effort, and posture during movement. \\
\midrule
\multirow{9}{=}{\centering\textbf{Spatial Context}}  & Atrium & Large interior space characterised by openness, verticality, and daylight penetration. \\
& Corridor & Linear passage organising horizontal movement between programmatic zones. \\
& Staircase & Vertical element structuring movement between levels. \\
& Threshold & Transitional zone marking a spatial or environmental shift. \\
& Entrance & Arrival space mediating exterior--interior transition and orientation. \\
& Workstation & Task-oriented space for focused work. \\
& Alcove & Partially enclosed recess enabling visual and acoustic separation. \\
& Lounge & Informal gathering area supporting pausing and social occupation. \\
& Outdoor-Adjacent & Edge condition where exterior environmental qualities permeate the interior. \\
\midrule
\multirow{5}{=}{\centering\textbf{Interaction Type}} & Attending & Experience shaped through receptive presence and sensory exposure without overt action. \\
& Traversing & Experience unfolds through bodily movement across space rather than manipulation. \\
& Handling & Physical engagement through touching, moving, or manipulating environmental elements. \\
& Adjusting-To & Reciprocal coupling where environmental behaviour shifts in response to presence or movement. \\
& Coordinating-With & Interaction emerges through mutual adjustment and co-presence of multiple occupants. \\
\bottomrule
\end{tabular}
\end{table*}

\begin{table*}[!ht]
\caption{Layer~1 codebook used for mapping the biophilic design space (Part 2).}
\Description{Codebook for Layer 1 showing dimensions, codes, and definitions for Affective Aim, Natural Reference, Temporal Behaviour, and Level of Mediation.}
\label{tab:layer1-codebook-continued}
\centering
\begin{tabular}{@{}
  >{\centering\arraybackslash}p{0.18\textwidth}
  p{0.20\textwidth}
  p{0.58\textwidth}
@{}}
\toprule
\textbf{Dimension} & \textbf{Code} & \textbf{Definition} \\
\midrule
\multirow{10}{*}{\textbf{Affective Aim}} & Calm / Relaxation & Ease and reduced arousal (e.g., soothing, soft, quiet). \\
& Mental Reset & Refreshed or cleared mental state through transition or embodied focus. \\
& Curiosity / Gentle Fascination & Effortless interest drawn by subtle variation or sensory intrigue. \\
& Playfulness / Light Delight & Lighthearted enjoyment or bodily fun through optional exploration. \\
& Connection / Empathy & Emotional resonance or relational closeness to nature or place. \\
& Natural Vitality / Aliveness & Sense of liveliness conveyed through movement or dynamic processes. \\
& Authenticity / Groundedness & Anchoring through raw, honest, or unfiltered materials and cues. \\
& Comfort / Safety / Reassurance & Physical or psychological ease and protection. \\
& Social Warmth / Togetherness & Positive social presence and ease in shared environments. \\
& Neutral & No affective orientation articulated. \\
\midrule
\multirow{9}{*}{\textbf{Natural Reference}} & Light \& Shadow & Sunlight, reflection, filtering, shimmer, shadow-play, diurnal cues. \\
& Wind / Airflow & Breezes, drafts, pressure changes, ventilation-like sensations. \\
& Water & Flowing, dripping, rippling, wetness, or wave-like motion. \\
& Vegetation & Trees, leaves, greenery, branches, plant density. \\
& Terrain / Ground Texture & Underfoot variation such as uneven ground, logs, roots, stepping surfaces. \\
& Animal Life & Cues from animals (e.g., birdsong, rustling, fauna-like sounds). \\
& Seasonality / Weather & Temperature, humidity, atmospheric pressure, seasonal change, natural odours. \\
& Open Landscape & Views, distance, horizon, spatial depth, openness. \\
& Natural Randomness & Irregularity, disorder, emergent or non-designed spatial logic. \\
\midrule
\multirow{7}{*}{\textbf{Temporal Behaviour}} & Static & No perceptible change over time. \\
& Perpetual & Continuous, steady-state behaviour without clear beginning or end. \\
& Gradual & Slow, subtle change unfolding over time. \\
& Cyclical / Rhythmic & Predictable recurring temporal patterns. \\
& Intermittent / Sporadic & Irregular, unpredictable bursts or fluctuations. \\
& Layered & Multiple temporalities occurring simultaneously. \\
& Ephemeral / Fleeting & Short-lived, transient sensory events. \\
\midrule
\multirow{4}{*}{\textbf{Level of Mediation}} & Natural & Experience arises with minimal human shaping. \\
& Architectural & Built form filters or structures natural qualities. \\
& Technological & Interactive or automated systems primarily produce or shape the experience. \\
& Hybrid / Mixed & Natural, architectural, and technological mediation jointly contribute. \\
\bottomrule
\end{tabular}
\end{table*}


